
\subsection{Motivation et Architecture Générale}

Nous adoptons une architecture
\textbf{Mixture of Experts (MoE)} inspirée des modèles d'apprentissage
automatique modernes. L'idée centrale est de diviser le problème en
sous-problèmes plus simples, chacun traité par un agent spécialisé
(\textit{expert}), et de confier à un \textbf{meta-agent} le soin de
sélectionner dynamiquement l'expert le plus adapté à la situation courante.

\subsection{Les Experts Spécialisés}

Trois agents experts sont entraînés indépendamment via l'algorithme
\textbf{PPO} (\textit{Proximal Policy Optimization}) de la bibliothèque
Ray RLlib. Chacun agit sur un sous-espace de l'action globale, représentée
par un vecteur à 6 dimensions $\mathbf{a} = [a_0, a_1, a_2, a_3, a_4, a_5]$
où chaque $a_i \in [0, 16]$.

\subsubsection{ScalingExpert}

Le \textit{ScalingExpert} contrôle le \textbf{scaling horizontal} des
fonctions, c'est-à-dire le nombre de réplicas déployés. Il agit sur les
dimensions $a_0$ et $a_1$ (les deux premières fonctions), en fixant les
autres à 8 (valeur neutre).

Sa fonction de récompense pénalise uniquement la latence :
\[
r_{\text{scaling}} = -15{,}0 \times \text{latence\_norm}
\]

\subsubsection{PowerExpert}

Le \textit{PowerExpert} gère l'\textbf{alimentation des nœuds} en décidant
d'allumer ou d'éteindre des nœuds du cluster. Il agit sur les dimensions
$a_2$ et $a_3$, avec une contrainte de sécurité garantissant qu'au moins
10\% des nœuds restent actifs.

Sa fonction de récompense équilibre économie d'énergie et latence :
\[
r_{\text{power}} = 10{,}0 \times \frac{n_{\text{off}}}{n_{\text{total}}}
                 - 15{,}0 \times \text{latence\_norm}
\]

\subsubsection{SchedulingExpert}

Le \textit{SchedulingExpert} choisit la \textbf{stratégie de placement}
des conteneurs : consolidation (remplir les nœuds) ou spreading (répartir
la charge). Il agit sur les dimensions $a_4$ et $a_5$.

Sa fonction de récompense favorise une utilisation CPU dans la plage optimale :
\[
r_{\text{sched}} = -5{,}0 \times \text{latence\_norm}
+ \begin{cases} +5{,}0 & \text{si } 0{,}4 < \text{CPU} < 0{,}8 \\ -2{,}0 & \text{sinon} \end{cases}
\]

\subsubsection{Résumé de l'espace d'action}

\begin{table}[h]
\centering
\begin{tabular}{|l|c|c|c|c|c|c|}
\hline
\textbf{Expert} & $a_0$ & $a_1$ & $a_2$ & $a_3$ & $a_4$ & $a_5$ \\
\hline
ScalingExpert   & VAR & VAR & 8 & 8 & 8 & 8 \\
PowerExpert     & 8 & 8 & VAR & VAR & 8 & 8 \\
SchedulingExpert & 8 & 8 & 8 & 8 & VAR & VAR \\
\hline
\end{tabular}
\caption{Répartition de l'espace d'action entre les trois experts
(8 = valeur neutre, VAR = dimension contrôlée par l'expert).}
\label{tab:action_space}
\end{table}

\subsection{Le Meta-Agent (Gating Network)}

\subsubsection{Rôle et Fonctionnement}

Le meta-agent joue le rôle de \textbf{gating network} : à chaque pas de
temps, il observe l'état global du cluster et décide quel expert doit agir.
L'observation est un vecteur de 8 dimensions :
\[
\mathbf{o} = [\text{latence\_norm},\ \text{CPU\_util},\ r_1,\ r_2,\ r_3,\ r_4,\ r_5,\ r_6]
\]
où $r_i$ est le nombre de réplicas de la fonction $i$.

\subsubsection{Entraînement séquentiel en deux phases}

L'entraînement global suit un protocole \textbf{séquentiel en deux phases}, conçu pour garantir la stabilité du modèle final et éviter le problème de \textit{cible mouvante} (\textit{moving target problem}).

\textbf{Phase 1 --- Entraînement des experts.}
Les trois experts (ScalingExpert, PowerExpert, SchedulingExpert) sont entraînés \textbf{indépendamment} via l'algorithme PPO jusqu'à convergence. Chaque expert est entraîné pendant environ 120 itérations ($\approx 30$ minutes). Une fois cette phase terminée, les poids de chaque expert sont \textbf{gelés} (\textit{frozen}) : ils ne seront plus mis à jour lors de la phase suivante.

\textbf{Phase 2 --- Entraînement du meta-agent.}
Le meta-agent est entraîné sur les experts figés de la Phase~1. Cette contrainte est essentielle : si les experts continuaient à évoluer pendant l'entraînement du meta-agent, ce dernier apprendrait à composer avec des politiques instables, rendant la convergence impossible. En fixant les experts au préalable, on garantit que le meta-agent apprend à sélectionner parmi des politiques stables et prévisibles.

Le meta-agent est entraîné sur 1\,000 itérations PPO avec un taux d'apprentissage $\eta = 10^{-4}$ et un coefficient d'entropie de $0{,}05$. Sa fonction de récompense unifiée est :
\[
r = -10{,}0 \times \text{latence\_norm} + 25{,}0 \times \frac{n_{\text{off}}}{n_{\text{total}}}
\]

Ce protocole en deux phases séquentielles, en entraînant d'abord les experts puis le meta-agent sur ces experts figés, assure la stabilité du système et évite les interactions circulaires qui empêcheraient la convergence.

\subsection{Discussion et Limites}

\subsubsection{Surpondération de l'énergie}

La fonction de récompense unifiée accorde un poids de $+25$ à l'énergie
contre $-10$ à la latence. Ce déséquilibre pousse le meta-agent à maximiser
l'extinction de nœuds même au détriment des SLAs. Une pondération
adaptative ou des contraintes dures sur la latence ($\text{latence} < 1500$ms)
permettraient d'obtenir un meilleur compromis.

\subsubsection{Entraînement séquentiel et non-optimalité}

Les experts étant figés avant l'entraînement du meta-agent, l'optimalité
jointe n'est pas garantie. Un entraînement conjoint (\textit{joint
training}) ou un mécanisme de \textit{fine-tuning} des experts en Phase 2
pourrait améliorer les performances globales.

\subsubsection{Absence de coordination inter-experts}

Les trois experts partagent la même observation mais n'ont aucun mécanisme
de communication entre eux. En pratique, ScalingExpert et PowerExpert ont
des objectifs contradictoires : l'un ajoute des réplicas (consomme plus),
l'autre éteint des nœuds (réduit les ressources disponibles). Une
architecture avec mémoire partagée ou communication explicite pourrait
réduire ces conflits.

La distribution des temps de latence illustrée sur la
Figure latency box confirme que certaines stratégies de référence
maintiennent de meilleures garanties de temps de réponse grâce à une gestion
plus fine et coordonnée des ressources, ce que l'architecture MoE actuelle
ne peut pas encore égaler en l'absence de coordination inter-experts.
Un mécanisme de \textit{multi-objective RL} ou une mémoire partagée entre
experts permettrait de dépasser cette limitation.


\subsubsection{Généralisation}

La charge de travail est tirée aléatoirement dans $[10, 45]$ RPS à chaque
épisode. Le meta-agent apprend donc une politique moyennée sur cet
ensemble de charges, ce qui peut le rendre sous-optimal pour des régimes
de charge spécifiques. Un curriculum d'entraînement progressif améliorerait
la robustesse.
